\sscp{Evidence of contact-induced change}{02-03-01}
\ili{As} shown in the following chapters,
under the surface it appears that historically there were multiple
incursions by Austronesian-speaking populations into the region,
who came into contact with pre-existing populations speaking
non-Austronesian (Papuan) languages that were typologically predominantly SOV.
There was bilingualism and intermarriage over millennia,
which resulted in Papuan languages influencing not only the vocabulary
(which by itself can result from relatively superficial contact),
but also the structures
and semantics of Austronesian languages, and \it{vice versa}.
Having both island and mountain ecologies, both tropical rainforest
and semi-arid climates, and newly encountered flora
and fauna in the region all meant that patterns of contact
and diffusion are quite different from island group to island
group, regardless of whether the contact was with Austronesian or Papuan speaking populations.
For example different patterns of localised contact provide the
most likely explanation for \ili{Havu}
and \ili{Dhao} being relatively closely related in their historical
phonologies and parts of their lexicon,
yet very different in their syntax and typologies \citep{Grimes2010a}.
The local contact histories in northern \ili{Buru},
northern Seram and northern Flores are significantly different
from those in southern \ili{Buru},
southern Seram and southern Flores respectively
\citep[11]{Grimes2000a,GrimesTherikGrimesJacob1997}.

\citet{ThomasonKaufman1988,Schapper2015,Klamer2019}, and
\citet{GrimesCIP} point out that the kinds of contact-induced
structural and typological changes found in languages of the region
and discussed in this section are best explained as reflecting
intense contact over a long period of time.
\citet[14--16]{Klamer2019} walks carefully through the logic of how
linguistic features that are not found in Austronesian languages
in the west and north, but are commonly found in both the Papuan languages
and the Austronesian languages of our target region can be attributed
to contact with the Papuan languages.
It is particularly important to understand the significance of
Austronesian languages that are typologically SVO shifting away
from structures that are commonly associated this typology,
and taking on structures that are commonly associated
with SOV typologies, but sometimes are only rarely associated with SVO typologies.
This is highly significant for understanding the linguistic prehistory of the region.

There are a variety of typological and lexical features in the languages
\ili{east} of the Blust line that are not typically found in languages to
its west and its north. These features tend to have a scattered
and/or non-uniform distribution meaning they cannot be attributed
to a single proto-language, or often even to multiple proto-languages.
Instead, these features appear to have been acquired by contact
after the break-up of proto-languages.

So for instance, the preposed possessor introduced in \srf{01-05}
is the normal pattern throughout our target region.
Yet within the \tsc{\ili{Bima}-Lembata} subgroup both preposed
and postposed possessors are found, with the preposed possessor
the most common order within the \tsc{Flores-Lembata} node,
yet even within \tsc{Flores-Lembata} post possessors also occur.

Similarly, in the lexicon \citet{DenhamDonohue2009,DonohueDenham2009,Schapper2015}
have pointed out that the pseudo-etymon {\#}muku `banana'
has a Papuan source and is frequent in the southern region
of our target region. But it co-occurs with reflexes retaining PMP *punti `banana'
(as well as other diverse etyma) even among members within the same node of a subgroup.
So, for instance, the \tsc{Timor-Babar} language \ili{Tetun} has
*punti > \it{hudi} while its sister \ili{Habun} within the \tsc{Eastern Timor}
node has {\#}muku > \it{muku}. Similarly, in the \tsc{Southwest Maluku}
node of \tsc{Timor-Babar} reflexes of *punti predominate,
but \ili{Kisar} has \it{muʔu}, while reflexes of **biak occur in
three other nodes of \tsc{Southwest Maluku}.
See \srf{02-03-01-07}, \frf{15-11}, and appendix \ref{ch:BanTer}
for more details.

Attempting to explain such distributions as the result of inheritance
from a common ancestor is futile. Instead, different subgroups, nodes,
and languages have each undergone contact in different ways and to different
extents which have resulted in the kinds of uneven distribution that we see.
A selection of the kinds of features that are found in such an uneven manner
is given in \trf{02-Features} divided into typological and lexical features.
This list is not intended to be exhaustive.
Some of these features are due to contact with Papuan languages,
while others may be due to contact between Austronesian
languages of the region after their arrival, and some may be due to both.
We discuss these features in more detail in chapter \ref{ch:MorLes}.

\begin{table}
	\centering\tcp{Contact features Austronesian languages of Wallacea}{02-Features}
	\st{6pt}\begin{tabular}{l@{ }ll}\lsptoprule
 \mc{2}{l}{Feature} & see \\	\midrule
 \mc{2}{l}{Preposed possessor} & \srf{02-03-01-01} \\
 \mc{2}{l}{Alienable-inalienable distinction} & \srf{02-03-01-01} \\
 \mc{2}{l}{Clause-final negation} & \srf{02-03-01-02} \\
 \mc{2}{l}{Postpositions} & \srf{02-03-01-03} \\
 \mc{2}{l}{Complex and non-decimal numerals} & \srf{02-03-01-04} \\
 \mc{2}{l}{absence of velar nasal /ŋ/} & \srf{13-01-01} \\
 \mc{2}{l}{subject prefixes on verbs} & \srf{13-03-03} \\
 \mc{2}{l}{Human/non-human distinction} & \srf{07-04-03} \\	\midrule
 *kəRaŋ & `shell, turtle' & \srf{13-01-05-03} \\
 **mantəR & `cuscus' & chapter \ref{ch:Mar} \\
 {\#}muku & `banana' & \srf{02-03-01-07} \\
 `head leaf' & = `hair' & \srf{02-03-01-05} \\
 `inside' & = `seat of emotions, character' & \srf{02-03-01-06} \\
 *duRi `thorn' &> *duRi `bone' & \srf{15-04-02-03} \\
 \mc{2}{l}{kinship systems and terminologies} & \srf{07-02-03}, \srf{15-04-02-04}, \srf{15-04-03-02} \\
 **malip & `laugh, smile' & \srf{15-04-04-01} \\
 **maya & `tongue' & \srf{15-04-04-02} \\
		\lspbottomrule
	\end{tabular}
\end{table}

The presence of a cluster of the features in \trf{02-Features}
in sets of languages adds to the complexity of unravelling what
is due to retention, shared innovation, parallel development,
and/or what is due to contact.

Although our discussion of contact in this work focusses on the influence of Papuan
languages on Austronesian languages in the region,
it has also been shown that Austronesian languages influenced
Papuan languages \citep{KlamerReesinkVanStaden2008,ReesinkDunn2018}.
For example, the development of an inclusive-exclusive distinction
in the pronouns of some Papuan languages is attributed to Austronesian contact.
And some see the shift to SVO typology in some nodes of Papuan
language families that have SOV typology in other nodes as due
to contact with Austronesian languages \citep[190--193]{Voorhoeve1988}.
The contact reflected by such changes is not trivial,
and imply far more significant contact than is required for
simply borrowing a few words.
